%ClassifierSemantics.tex


%\chapter{Classifiers}

\comment{Note:  This Classifiers Semantics subclause is presented after Features, although it retains the paragraph numbering.}

\ksec{Classifiers Semantics (KerML 8.4.3.3)}

\pn \texttt{Classifiers} are \texttt{Types} that classify things in the modeled system, as distinct from features, which model the relations between them.  
\texttt{Classifiers} have no temporal semantics.\footnote{Temporal semantics have been moved from Classifier to Class at the behest of 
\bemph{Vince Moln\'ar}.}
Let $V_C$ be the set of all \texttt{Classifiers} in the design $D$:  $V_C \subset D$, then for all elements in the design, an element being defined a \texttt{Classifier} implies the element is a member of $V_C$.

\begin{restatable}{definition}{dfclassifier}~\blTitle{df-classifier}~  Define \texttt{classifier}.
\begin{align}
\vdash \forall c \in D~(\texttt{classifier}~c~ \rightarrow ~c \in V_C)~\label{eq:df-classifier}\tag{df-classifier}
\end{align}
\end{restatable}

\begin{restatable}{leanaxiom}{VC}~\leanTitle{VC}~ Define the set of \texttt{Classifiers} in Lean. Exact
parallel of \texttt{VT}/\texttt{VF} above, one more \texttt{ElementKind} tag down.

\texttt{\textcolor{leancolor}{axiom} \textcolor{leanyellow}{VC} : Set Element}
\label{VC}
\end{restatable}

\begin{restatable}{leanaxiom}{dfClassifierLean}~\leanTitle{df\_classifier}~ Define \texttt{classifier} in Lean.
\texttt{Classifier} already has a slot in \texttt{ElementKind}'s fourteen tags, so no new \texttt{DesignKind}
case is needed --- exact parallel of \texttt{df\_types}/\texttt{df\_feature} above.

\texttt{\textcolor{leancolor}{axiom} \textcolor{leanyellow}{df\_classifier} : $\forall$ e : Element,}

\texttt{~~($\exists$ C : Set Element, (DesignKind.ofElementKind ElementKind.Classifier, e, C) $\in$ Design) $\to$ e $\in$ VC}
\label{df_classifier}
\end{restatable}

\pn \texttt{Classifiers} are also \texttt{Types}.

\begin{restatable}{definition}{dfclassifier2} ~\blTitle{df-classifier2}~  \texttt{Classifiers} are also \texttt{Types}.
\begin{align}
\vdash V_C~\subset~V_T~\label{eq:df-classifier2}\tag{df-classifier2}
\end{align}
\end{restatable}

\begin{restatable}{leanaxiom}{dfClassifier2Lean}~\leanTitle{df\_classifier2}~ \texttt{Classifiers} are also
\texttt{Types}, in Lean. Not derivable from \texttt{VC}'s/\texttt{VT}'s own definitions (each is
independently primitive, same situation \texttt{df\_types}/\texttt{df\_feature} were already in), so this
stays an axiom here too.

\texttt{\textcolor{leancolor}{axiom} \textcolor{leanyellow}{df\_classifier2} : VC $\subseteq$ VT}
\label{df_classifier2}
\end{restatable}

\pn A \texttt{Classifier} is just the same as a \texttt{Type}, but it also is a member of $V_C$.


%\ksec{Value Semantics}
%
%\pn A \texttt{Value}\index{value}, $v$, is an Element\index{Element} of a \texttt{Type}, $T$.  A \texttt{Type}\index{type} is a set of \texttt{Values}.  The value of a Value is a ZFC set.
%
%\begin{restatable}{definition} ~\blTitle{df-value}~  \texttt{Values} are elements of a \texttt{Type}.
%\begin{align}
%T~\in~V_T~\wedge~v~\in~T  \label{eq:df-value}\tag{df-value}
%\end{align}
%\end{restatable}
%
%\pn A \texttt{Value} may be expressed as a literal, a \texttt{Function}\index{function}, an \texttt{Expression}\index{expression}, or a \texttt{Predicate}\index{predicate}.
%
%\pn A \texttt{DataValue}, $dv$, is pure information so is not a \texttt{Classifier} or a \texttt{Feature} (things modeled in the system).  Let $DV$ be the set of all \texttt{DataValues}.
%
%\begin{restatable}{definition} ~\blTitle{df-datavalue}~  \texttt{DataValues} are elements of a \texttt{Type} but not a \texttt{Classifier} or a \texttt{Feature}.
%\begin{align}
%DV~\subset~(V_T~\backslash~(V_C~\cup~V_F~)~)~\wedge~dv~\in~DV  \label{eq:df-datavalue}\tag{df-datavalue}
%\end{align}
%\end{restatable}
%
